%% series_2_ordinarios.tex -- version ampliada de series_2 (soluciones en serie en puntos ordinarios)
%% Agregados: recta con puntos singulares, tabla de coeficientes, diagrama del metodo,
%% tabla de errores tipicos, comparaciones grafico-numericas (Hermite, Ej. 3, R=1, Airy),
%% tabla resumen y ejercicios con respuesta. En todo el archivo la forma estandar es
%% y''+P(x)y'+Q(x)y=0 (mayusculas, como en la seccion de puntos singulares).
%% Requiere en el preambulo del documento principal:
%%   \usepackage{pgfplots}  \pgfplotsset{compat=1.18}
%% Convenciones de notacion del capitulo: en la forma estandar y''+P(x)y'+Q(x)y=0
%% los coeficientes se llaman P y Q; en la forma de Frobenius
%% x^2 y''+x p(x) y'+q(x) y=0 se usan p=(x-x0)P y q=(x-x0)^2 Q (minusculas).
 
%% --- Correccion de \ref a cajas numeradas (ver sistemas_2.tex) --------------
\makeatletter
\ifdefined\bbereffixdone\else
  \def\bbereffixdone{}
  \tcbset{before upper={%
    \ifx\thetcbcounter\@empty\else
      \protected@edef\@currentlabel{\thetcbcounter}%
      \ifdefined\phantomsection\phantomsection\fi
    \fi}}
\fi
\makeatother
 
%% --- Estilo comun de los graficos de este capitulo --------------------------
\pgfplotsset{
  bbeserie/.style={width=.94\linewidth, height=6.2cm, axis lines=middle,
    xlabel=$x$, ylabel=$y$, tick label style={font=\small}, label style={font=\small},
    legend style={font=\footnotesize, draw=bbe!40, fill=white, fill opacity=.9, text opacity=1},
    grid=major, grid style={draw=bbe!12}, clip=true,
    every axis plot/.append style={thick}},
  bbeosc/.style={bbeserie, axis lines=left}
}
\providecommand{\bbezero}[2]{\addplot[bbegray!70, thin, forget plot] coordinates {(#1,0) (#2,0)};}
\tikzset{
  bbecaja/.style={draw=bbe, fill=bbe!8, rounded corners=2pt, text=black, align=center,
    font=\small\sffamily, inner sep=4pt, text width=3.6cm, minimum height=1.1cm},
  bbeflujo/.style={-{Stealth[length=5pt]}, thick}
}
 
 
Ahora consideraremos la ecuación diferencial lineal homogénea de segundo orden escrita en forma estándar:
\[ y'' + P(x)y' + Q(x)y = 0. \]
 
Nuestro objetivo es determinar una representación en serie de la solución general a esta ecuación diferencial centrada en $x = x_0$. Veremos que la existencia y forma de la solución dependen del comportamiento de las funciones $P$ y $Q$ en $x = x_0$.
 
\begin{definition}
    El punto $x = x_0$ se llama un \textbf{punto ordinario} de la ecuación diferencial
\begin{equation}
y'' + P(x)y' + Q(x)y = 0,
\label{eq:s2-forma-estandar}
\end{equation}
si $P$ y $Q$ son \textit{ambas} analíticas en $x = x_0$. Cualquier punto que no sea un punto ordinario de \eqref{eq:s2-forma-estandar} se llama un \textbf{punto singular} de la ecuación diferencial.
\end{definition}
 
\begin{example}
    La ecuación diferencial
\[ y'' + \frac{1}{x^2 - 4}y' + \frac{1}{x + 1}y = 0 \]
tiene
\[ P(x) = \frac{1}{x^2 - 4} \quad \text{y} \quad Q(x) = \frac{1}{x + 1}. \]
 
Vemos por inspección que los únicos puntos en los que $P$ falla en ser analítica son $x = \pm 2$, mientras que $Q$ es analítica en todos los puntos excepto en $x = -1$. En consecuencia, los únicos puntos singulares de la ecuación diferencial son $x = -1, \pm 2$ (Figura~\ref{fig:s2-singulares}). Todos los demás puntos son puntos ordinarios. \hfill $\square$
\end{example}
 
\begin{figure}[H]
\centering
\begin{tikzpicture}[x=1.7cm]
  \draw[->, bbegray] (-3.6,0) -- (3.7,0) node[right, text=black, font=\small] {$x$};
  \foreach \t in {-3,...,3} \draw[bbegray] (\t,-.07) -- (\t,.07) node[below=3pt, text=black, font=\footnotesize] {$\t$};
  \draw[bbegreen, ultra thick] (-3.5,0) -- (-2.08,0);
  \draw[bbegreen, ultra thick] (-1.92,0) -- (-1.08,0);
  \draw[bbegreen, ultra thick] (-.92,0) -- (1.92,0);
  \draw[bbegreen, ultra thick] (2.08,0) -- (3.55,0);
  \foreach \s in {-2,-1,2} \node[bbered, font=\Large] at (\s,0) {$\times$};
  \node[text=black, font=\footnotesize, above=3pt] at (-2,0) {$P$ falla};
  \node[text=black, font=\footnotesize, above=3pt] at (-1,0) {$Q$ falla};
  \node[text=black, font=\footnotesize, above=3pt] at (2,0) {$P$ falla};
  \node[bbegreen, font=\footnotesize\sffamily] at (0.5,.45) {puntos ordinarios};
\end{tikzpicture}
\caption{Puntos singulares ($\times$) y puntos ordinarios (verde) de la ecuación del ejemplo.}
\label{fig:s2-singulares}
\end{figure}
 
\begin{remark}
    En esta sección restringimos nuestra atención a los puntos ordinarios. Dado que las funciones $P$ y $Q$ son analíticas en un punto ordinario, podríamos sospechar que cualquier solución de \eqref{eq:s2-forma-estandar} válida en $x = x_0$ también es analítica allí y, por lo tanto, se puede representar como una serie de potencias convergente
\[ y(x) = \sum_{n=0}^{\infty} a_n(x - x_0)^n, \]
para constantes apropiadas $a_n$. Este es de hecho el caso, pero antes de enunciar el resultado general, consideramos un ejemplo para ilustrar la idea anterior.
\end{remark}
 
\begin{example}\label{ex:s2-cos}
    Determine dos soluciones en serie de potencias linealmente independientes para la ecuación diferencial
\begin{equation}
y'' + y = 0,
\label{eq:s2-oscilador}
\end{equation}
centrada en $x = 0$. Identifique las soluciones en términos de funciones elementales familiares.
 
\noindent \textbf{Solución:} \quad Dado que $x = 0$ es un punto ordinario de la ecuación diferencial, intentamos una solución en serie de potencias de la forma
\begin{equation}
y(x) = \sum_{n=0}^{\infty} a_n x^n.
\label{eq:s2-ansatz}
\end{equation}
 
Procedemos de manera similar al método de coeficientes indeterminados sustituyendo \eqref{eq:s2-ansatz} en \eqref{eq:s2-oscilador} y determinando los valores de los $a_n$ tales que \eqref{eq:s2-ansatz} sea efectivamente una solución. Derivar \eqref{eq:s2-ansatz} dos veces con respecto a $x$ produce
\[ y'(x) = \sum_{n=1}^{\infty} n a_n x^{n-1}, \qquad y''(x) = \sum_{n=2}^{\infty} n(n-1)a_n x^{n-2}, \]
donde hemos desplazado el punto de inicio en las sumatorias sin pérdida de generalidad. Sustituyendo en \eqref{eq:s2-oscilador}, se deduce que \eqref{eq:s2-ansatz} define una solución, siempre que
\[ \sum_{n=2}^{\infty} n(n-1)a_n x^{n-2} + \sum_{n=0}^{\infty} a_n x^n = 0. \]
 
Si reemplazamos $n$ por $k + 2$ en la primera sumatoria, y $n$ por $k$ en la segunda, el resultado es
\[ \sum_{k=0}^{\infty} (k+2)(k+1)a_{k+2}x^k + \sum_{k=0}^{\infty} a_k x^k = 0, \]
y combinando,
\[ \sum_{k=0}^{\infty} [(k+2)(k+1)a_{k+2} + a_k]x^k = 0. \]
 
Esto implica que los coeficientes de $x^k$ deben anularse para $k = 0, 1, 2, \dots$. En consecuencia, obtenemos la \textbf{relación de recurrencia}
\begin{equation}
a_{k+2} = -\frac{1}{(k+2)(k+1)}a_k, \qquad k = 0, 1, 2, \dots
\label{eq:s2-rec-cos}
\end{equation}
(dividimos por $(k+2)(k+1)$, que nunca es cero). La recurrencia \emph{salta de dos en dos}: los coeficientes pares dependen solo de $a_0$ y los impares solo de $a_1$, y ambos son libres. La tabla siguiente muestra el patrón:
\begin{center}
\small
\begin{tabular}{|c|c|c|c|c|c|c|}
\hline
\rowcolor{bbe!15}$k$ & $0$ & $1$ & $2$ & $3$ & $4$ & $5$ \\ \hline
$a_{k+2}$ & $-\dfrac{a_0}{2!}$ & $-\dfrac{a_1}{3!}$ & $+\dfrac{a_0}{4!}$ & $+\dfrac{a_1}{5!}$ & $-\dfrac{a_0}{6!}$ & $-\dfrac{a_1}{7!}$ \\ \hline
\end{tabular}
\end{center}
Reconocemos así los términos generales
\begin{equation}
\boxed{ a_{2n} = \frac{(-1)^n}{(2n)!}a_0, \qquad a_{2n+1} = \frac{(-1)^n}{(2n+1)!}a_1. }
\label{eq:s2-coef-cos}
\end{equation}
 
Por lo tanto, para todos los valores de $a_0, a_1$, una solución a la ecuación diferencial dada es
\begin{equation}
y(x) = a_0 \sum_{n=0}^{\infty} \frac{(-1)^n}{(2n)!}x^{2n} + a_1 \sum_{n=0}^{\infty} \frac{(-1)^n}{(2n+1)!}x^{2n+1}.
\label{eq:s2-sol-cos}
\end{equation}
 
Establecer $a_0 = 1$ y $a_1 = 0$ produce $y_1(x) = \sum_{n\ge 0} \frac{(-1)^n}{(2n)!}x^{2n}$, mientras que $a_0 = 0$ y $a_1 = 1$ produce $y_2(x) = \sum_{n\ge 0} \frac{(-1)^n}{(2n+1)!}x^{2n+1}$.
 
Aplicando el criterio de la razón (a la variable $u=x^2$, ver la nota de la sección anterior), ambas series convergen para todo $x$ real. Dado que $y_1$ y $y_2$ no son proporcionales, son linealmente independientes en cualquier intervalo. Se deduce que \eqref{eq:s2-sol-cos} es la solución general. De hecho, $y_1$ es la serie de Maclaurin de $\cos x$ y $y_2$ la de $\sin x$ (Tabla~\ref{tab:s1-maclaurin}). Por lo tanto,
\[ y(x) = a_0 \cos x + a_1 \sin x. \tag*{$\square$} \]
\end{example}
 
El desarrollo y la solución del problema anterior se basó esencialmente en los pasos que resume la Figura~\ref{fig:s2-metodo}; la Tabla~\ref{tab:s2-errores} reúne los tropiezos más frecuentes.
 
\begin{figure}[H]
\centering
\begin{tikzpicture}[node distance=.55cm, bbecaja/.append style={text width=3.3cm, minimum height=1.5cm}]
  \node[bbecaja] (p1) {\textbf{1.} Proponer\\ $y=\sum a_nx^n$};
  \node[bbecaja, right=of p1] (p2) {\textbf{2.} Derivar y sustituir en la ecuación};
  \node[bbecaja, right=of p2] (p3) {\textbf{3.} Corrimiento de índices: potencia común $x^k$};
  \node[bbecaja, below=.9cm of p3] (p4) {\textbf{4.} Igualar a cero cada coeficiente: \emph{recurrencia}};
  \node[bbecaja, left=of p4] (p5) {\textbf{5.} Resolver: patrón general o primeros términos};
  \node[bbecaja, left=of p5] (p6) {\textbf{6.} Radio $R$ e independencia lineal};
  \draw[bbeflujo] (p1) -- (p2); \draw[bbeflujo] (p2) -- (p3);
  \draw[bbeflujo] (p3) -- (p4); \draw[bbeflujo] (p4) -- (p5); \draw[bbeflujo] (p5) -- (p6);
\end{tikzpicture}
\caption{Método de series de potencias en un punto ordinario.}
\label{fig:s2-metodo}
\end{figure}
 
\begin{table}[H]
\centering
\small
\caption{Errores típicos al aplicar el método y cómo evitarlos.}
\label{tab:s2-errores}
\begin{tabular}{|p{5.6cm}|p{9.2cm}|}
\hline
\rowcolor{bbe!15}\textbf{Error típico} & \textbf{Cómo evitarlo} \\ \hline
Olvidar que $y'$ y $y''$ empiezan en $n=1$ y $n=2$ & Los términos que se anulan al derivar no aportan; comenzar la suma en el primer índice que no se anula. \\ \hline
Igualar coeficientes con potencias distintas & Antes de sumar, llevar \emph{todas} las series a $x^k$ (Tabla~\ref{tab:s1-corrimiento}). \\ \hline
Perder los términos "sueltos" de $k=0,1$ & Si las sumas empiezan en índices distintos, sacar a mano los primeros términos (ver el Ejemplo~\ref{ex:s2-tres}). \\ \hline
Dividir por una expresión que puede ser $0$ & Verificar que el factor de $a_{k+2}$, $(k+2)(k+1)$, nunca se anula para $k\ge 0$. \\ \hline
Creer que hay una sola constante libre & Siempre hay \textbf{dos}: $a_0$ y $a_1$ (los valores $y(x_0)$ e $y'(x_0)$). \\ \hline
\end{tabular}
\end{table}
 
El siguiente teorema justifica los pasos anteriores y muestra que la técnica se puede aplicar alrededor de cualquier punto ordinario de una ecuación diferencial.
 
\begin{theorem}\label{thm:s2-ordinario}
    Sean $P$ y $Q$ analíticas en $x = x_0$, y suponga que sus expansiones en serie de potencias son válidas para $|x - x_0| < R$. Entonces la solución general a la ecuación diferencial
\begin{equation}
y'' + P(x)y' + Q(x)y = 0
\label{eq:s2-teorema}
\end{equation}
puede representarse como una serie de potencias centrada en $x = x_0$, con un radio de convergencia de \textit{al menos} $R$. Los coeficientes en esta solución en serie pueden determinarse en términos de $a_0$ y $a_1$ sustituyendo directamente
\[ y(x) = \sum_{n=0}^{\infty} a_n(x - x_0)^n \]
en \eqref{eq:s2-teorema}. La solución resultante es de la forma
\[ y(x) = a_0 y_1(x) + a_1 y_2(x), \]
donde $y_1$ y $y_2$ son soluciones linealmente independientes de \eqref{eq:s2-teorema} en el intervalo de existencia. Si se imponen las condiciones iniciales $y(x_0) = \alpha$, $y'(x_0) = \beta$, entonces $a_0 = \alpha, a_1 = \beta$.
\end{theorem}
 
\begin{remark}
\textbf{¿Qué $R$ garantiza el teorema?} Si $P$ y $Q$ son funciones racionales, la serie de $P$ y $Q$ centrada en $x_0$ converge hasta la singularidad más cercana, \emph{real o compleja} (Teorema del radio de una función racional, en el repaso de series). Por lo tanto, el radio de la solución es \textbf{al menos} la distancia de $x_0$ al punto singular más cercano de la ecuación. Si $P$ y $Q$ son polinomios, $R=\infty$.
\end{remark}
 
A continuación algunos ejemplos para darle uso al teorema anterior.
 
\begin{example}\label{ex:s2-hermite}
  Determine dos soluciones en serie linealmente independientes en potencias de $x$ para
\begin{equation}
y'' - 2xy' - 4y = 0,
\label{eq:s2-hermite}
\end{equation}
y encuentre el radio de convergencia de estas soluciones.
 
\noindent \textbf{Solución:} \quad El punto $x = 0$ es un punto ordinario de la ecuación diferencial, y por lo tanto el Teorema~\ref{thm:s2-ordinario} se puede aplicar con $x_0 = 0$. En este caso, tenemos
\[ P(x) = -2x, \qquad Q(x) = -4. \]
 
Dado que ambos son polinomios, sus expansiones en serie de potencias alrededor de $x = 0$ son válidas para todo $x$, y por lo tanto la solución general de \eqref{eq:s2-hermite} puede expresarse en la forma $y(x) = \sum_{n=0}^{\infty} a_n x^n$, y esta solución en serie de potencias convergerá para todo $x$ real. Derivando y sustituyendo, obtenemos
\[ \sum_{n=2}^{\infty} n(n-1)a_n x^{n-2} - 2\sum_{n=1}^{\infty} n a_n x^n - 4\sum_{n=0}^{\infty} a_n x^n = 0. \]
 
Ahora redefinimos los rangos en la sumatoria para obtener un $x^k$ común en todos los términos. Esto se logra reemplazando $n$ con $k + 2$ en la primera sumatoria, y, por consistencia en la notación, reemplazamos $n$ con $k$ en las otras sumatorias. El resultado es
\[ \sum_{k=0}^{\infty} [(k+2)(k+1)a_{k+2} - 2ka_k - 4a_k]x^k = 0. \]
 
Esta ecuación requiere que el coeficiente de $x^k$ se anule, y por lo tanto, obtenemos la relación de recurrencia $(k+2)(k+1)a_{k+2} = 2(k+2)a_k$, es decir,
\begin{equation}
a_{k+2} = \frac{2}{k+1}a_k, \qquad k = 0, 1, 2, \dots
\label{eq:s2-rec-hermite}
\end{equation}
 
Vemos por esta relación que, como en el Ejemplo~\ref{ex:s2-cos}, todos los coeficientes pares pueden expresarse en términos de $a_0$, mientras que todos los coeficientes impares pueden expresarse en términos de $a_1$.
 
\vspace{0.2cm}
\noindent \textbf{Coeficientes pares} ($k=0,2,4,\dots$): $a_2 = 2a_0$, \; $a_4 = \frac{2}{3}a_2 = \frac{2^2}{3}a_0$, \; $a_6 = \frac{2}{5}a_4 = \frac{2^3}{1\cdot 3\cdot 5}a_0$, y en general
\[ a_{2n} = \frac{2^n}{1 \cdot 3 \cdot 5 \cdot \dots \cdot (2n - 1)}a_0. \]
 
\noindent \textbf{Coeficientes impares} ($k=1,3,5,\dots$): $a_3 = a_1$, \; $a_5 = \frac{2}{4}a_3 = \frac{1}{2!}a_1$, \; $a_7 = \frac{2}{6}a_5 = \frac{1}{3!}a_1$, \; $a_9 = \frac{2}{8}a_7 = \frac{1}{4!}a_1$, y en general
\[ a_{2n+1} = \frac{1}{n!}a_1. \]
 
Sustituyendo de nuevo, dos soluciones linealmente independientes de \eqref{eq:s2-hermite} en $(-\infty, \infty)$ son
\[ y_1(x) = 1 + \sum_{n=1}^{\infty} \frac{2^n}{1 \cdot 3 \cdot \dots \cdot (2n-1)}x^{2n}, \qquad y_2(x) = \sum_{n=0}^{\infty} \frac{1}{n!}x^{2n+1}. \]
 
\noindent\textbf{Identificación.} La serie $y_2=x\sum_{n\ge 0}\frac{(x^2)^n}{n!}$ es la de $x\,e^{x^2}$, de modo que $y_2(x)=x\,e^{x^2}$ (verifícalo: $y_2'=(1+2x^2)e^{x^2}$, $y_2''=(6x+4x^3)e^{x^2}$, y $y_2''-2xy_2'-4y_2=0$). Además, como $1\cdot3\cdots(2n-1)=\frac{(2n)!}{2^n n!}$, se puede escribir $y_1(x)=\sum_{n\ge 0}\frac{4^n\,n!}{(2n)!}\,x^{2n}$. \hfill $\square$
\end{example}
 
\begin{figure}[H]
\centering
\begin{tikzpicture}
\begin{axis}[bbeserie, height=6cm, xmin=0, xmax=1.7, ymin=0, ymax=30, restrict y to domain=-1:60,
  legend pos=north west, samples=120]
\addplot[black, very thick, domain=0:1.7] {x*exp(x^2)};
\addplot[bbered, domain=0:1.7] {x + x^3};
\addplot[bbeamber, domain=0:1.7] {x + x^3 + x^5/2};
\addplot[bbegreen, domain=0:1.7] {x + x^3 + x^5/2 + x^7/6};
\addplot[bbeblue, domain=0:1.7] {x + x^3 + x^5/2 + x^7/6 + x^9/24};
\legend{$xe^{x^2}$, hasta $x^3$, hasta $x^5$, hasta $x^7$, hasta $x^9$}
\end{axis}
\end{tikzpicture}
\caption{Ejemplo~\ref{ex:s2-hermite}: sumas parciales de $y_2=\sum x^{2n+1}/n!$ frente a $xe^{x^2}$. Cada término adicional alarga el tramo en que la aproximación es buena.}
\label{fig:s2-hermite}
\end{figure}
 
En general, no será posible encontrar la recurrencia entre los coeficientes de las respectivas series de potencias, o ésta será tan complicada que no se pueda resolver; es por esto que en la mayoría de las ocasiones debemos ser felices con obtener una cantidad finita de términos de cada solución en serie de potencias, por ejemplo:
 
\begin{example}\label{ex:s2-tres}
    Determine los términos hasta $x^5$ en cada una de las dos soluciones en serie de potencias linealmente independientes para
\[ y'' + (2 - 4x^2)y' - 8xy = 0 \]
centradas en $x = 0$. Encuentre también el radio de convergencia de estas soluciones.
 
\noindent \textbf{Solución:} \quad Las funciones $P(x) = 2 - 4x^2$ y $Q(x) = -8x$ son polinomios, y por lo tanto, a partir del Teorema~\ref{thm:s2-ordinario}, las soluciones en serie de potencias convergerán para todo $x$ real. Ahora determinamos las soluciones. Sustituir
\begin{equation}
y(x) = \sum_{n=0}^{\infty} a_n x^n
\label{eq:s2-ansatz-3}
\end{equation}
en la ecuación diferencial dada produce
\[ \sum_{n=2}^{\infty} n(n-1)a_n x^{n-2} + 2\sum_{n=1}^{\infty} n a_n x^{n-1} - 4\sum_{n=1}^{\infty} n a_n x^{n+1} - 8\sum_{n=0}^{\infty} a_n x^{n+1} = 0. \]
 
Reemplazando $n$ por $k + 2$ en la primera sumatoria, $k + 1$ en la segunda, y $k - 1$ en la tercera y cuarta, obtenemos
\[ \sum_{k=0}^{\infty} (k+2)(k+1)a_{k+2}x^k + 2\sum_{k=0}^{\infty} (k+1)a_{k+1}x^k - 4\sum_{k=2}^{\infty} (k-1)a_{k-1}x^k - 8\sum_{k=1}^{\infty} a_{k-1}x^k = 0. \]
Las sumas empiezan en índices distintos ($k=0,0,2,1$), así que separamos los términos $k = 0$ y $k = 1$ (las "sueltas" que advertía la Tabla~\ref{tab:s2-errores}):
\begin{align*}
(2a_2 + 2a_1) + (6a_3 + 4a_2 - 8a_0)x + \quad & \\
\sum_{k=2}^{\infty} \{(k+2)(k+1)a_{k+2} + 2(k+1)a_{k+1} - [4(k-1) + 8]a_{k-1}\}x^k &= 0.
\end{align*}
Igualar los coeficientes de todas las potencias $x^k$ a cero produce, para $k = 0$ y $k = 1$,
\begin{equation}
2a_2 + 2a_1 = 0 \quad \text{y} \quad 6a_3 + 4a_2 - 8a_0 = 0,
\label{eq:s2-inicio-3}
\end{equation}
y para $k \ge 2$,
\begin{equation}
(k+2)(k+1)a_{k+2} + 2(k+1)a_{k+1} - 4(k+1)a_{k-1} = 0,
\label{eq:s2-rec-3-previa}
\end{equation}
pues $4(k-1)+8=4(k+1)$. Se deduce de \eqref{eq:s2-inicio-3} que
\begin{equation}
a_2 = -a_1, \qquad a_3 = \frac{2}{3}(2a_0 + a_1),
\label{eq:s2-a2a3}
\end{equation}
y, dividiendo \eqref{eq:s2-rec-3-previa} por $k+1$, la relación de recurrencia general
\begin{equation}
a_{k+2} = \frac{4a_{k-1} - 2a_{k+1}}{k+2}, \qquad k = 2, 3, 4, \dots
\label{eq:s2-rec-3}
\end{equation}
(Observación: \eqref{eq:s2-rec-3} vale en realidad para todo $k\ge 0$ con la convención $a_{-1}=0$; las dos igualdades \eqref{eq:s2-inicio-3} son su caso $k=0,1$.)
 
La recurrencia involucra \emph{tres} coeficientes y es difícil de resolver en forma cerrada. Como solo se pidieron los términos hasta $x^5$, procedemos así. Ya tenemos $a_2$ y $a_3$ en términos de $a_0$ y $a_1$. Con $k = 2$ en \eqref{eq:s2-rec-3},
\[ a_4 = \frac{1}{4}(4a_1 - 2a_3) = \frac{1}{4}\left[ 4a_1 - \frac{4}{3}(2a_0 + a_1) \right] = \frac{2}{3}(a_1 - a_0), \]
y con $k = 3$,
\[ a_5 = \frac{1}{5}(4a_2 - 2a_4) = \frac{1}{5}\left[ -4a_1 - \frac{4}{3}(a_1 - a_0) \right] = \frac{4}{15}(a_0 - 4a_1). \]
Sustituyendo $a_2,\dots,a_5$ en \eqref{eq:s2-ansatz-3} y agrupando en $a_0$ y $a_1$,
\begin{align*}
y(x) &= a_0 \left( 1 + \frac{4}{3}x^3 - \frac{2}{3}x^4 + \frac{4}{15}x^5 + \dots \right) \\
&\quad + a_1 \left( x - x^2 + \frac{2}{3}x^3 + \frac{2}{3}x^4 - \frac{16}{15}x^5 + \dots \right).
\end{align*}
Así, dos soluciones linealmente independientes en $(-\infty, \infty)$ son
\[ y_1(x) = 1 + \frac{4}{3}x^3 - \frac{2}{3}x^4 + \frac{4}{15}x^5 + \dots, \qquad y_2(x) = x - x^2 + \frac{2}{3}x^3 + \frac{2}{3}x^4 - \frac{16}{15}x^5 + \dots \ . \tag*{$\square$} \]
\end{example}
 
\begin{remark}
\textbf{¿Qué tan buena es una solución truncada?} Aquí ocurre algo poco frecuente: la ecuación tiene una integral primera. Como
$\frac{d}{dx}\bigl[(2-4x^2)y\bigr]=(2-4x^2)y'-8xy$, la ecuación equivale a $\frac{d}{dx}\bigl[y'+(2-4x^2)y\bigr]=0$, es decir, $y'+(2-4x^2)y=C$. Con el factor integrante $e^{g(x)}$, $g(x)=2x-\frac43x^3$, se obtiene
\[ y(x)=e^{-g(x)}\Bigl[a_0+C\int_0^x e^{g(t)}\,dt\Bigr],\qquad C=a_1+2a_0. \]
Esto permite comparar con los polinomios de grado $5$ (Figura~\ref{fig:s2-ej3}): son excelentes cerca de $x=0$ y se separan de la solución exacta al alejarnos, a pesar de que $R=\infty$.
\end{remark}
 
\begin{figure}[H]
\centering
\begin{subfigure}[t]{.49\linewidth}
\centering
\begin{tikzpicture}
\begin{axis}[bbeserie, width=\linewidth, height=5.6cm, xmin=-1.5, xmax=1.5, ymin=-3, ymax=9,
  restrict y to domain=-10:20, legend pos=north west, legend style={font=\scriptsize}, samples=100]
\addplot[black, very thick, smooth] coordinates {(-1.5,-0.3530) (-1.45,-0.3750) (-1.4,-0.3926) (-1.35,-0.4035) (-1.3,-0.4053) (-1.25,-0.3958) (-1.2,-0.3730) (-1.15,-0.3357) (-1.1,-0.2834) (-1.05,-0.2165) (-1,-0.1362) (-0.95,-0.0445) (-0.9,0.0557) (-0.85,0.1613) (-0.8,0.2689) (-0.75,0.3752) (-0.7,0.4773) (-0.65,0.5726) (-0.6,0.6590) (-0.55,0.7353) (-0.5,0.8008) (-0.45,0.8553) (-0.4,0.8992) (-0.35,0.9333) (-0.3,0.9587) (-0.25,0.9765) (-0.2,0.9882) (-0.15,0.9952) (-0.1,0.9986) (-0.05,0.9998) (0,1.0000) (0.05,1.0002) (0.1,1.0013) (0.15,1.0042) (0.2,1.0097) (0.25,1.0187) (0.3,1.0317) (0.35,1.0497) (0.4,1.0735) (0.45,1.1041) (0.5,1.1425) (0.55,1.1903) (0.6,1.2492) (0.65,1.3212) (0.7,1.4092) (0.75,1.5169) (0.8,1.6490) (0.85,1.8117) (0.9,2.0135) (0.95,2.2658) (1,2.5844) (1.05,2.9913) (1.1,3.5176) (1.15,4.2081) (1.2,5.1283) (1.25,6.3756) (1.3,8.0981) (1.35,10.5242) (1.4,14.0150) (1.45,19.1525) (1.5,26.8976)};
\addplot[bbered, domain=-1.5:1.5] {1 + 4/3*x^3 - 2/3*x^4 + 4/15*x^5};
\legend{exacta, grado 5}
\end{axis}
\end{tikzpicture}
\caption{$y_1$ ($a_0=1,\ a_1=0$).}
\end{subfigure}\hfill
\begin{subfigure}[t]{.49\linewidth}
\centering
\begin{tikzpicture}
\begin{axis}[bbeserie, width=\linewidth, height=5.6cm, xmin=-1.5, xmax=1.5, ymin=-3, ymax=9,
  restrict y to domain=-10:20, legend pos=north west, legend style={font=\scriptsize}, samples=100]
\addplot[black, very thick, smooth] coordinates {(-1.5,-0.2881) (-1.45,-0.3435) (-1.4,-0.4082) (-1.35,-0.4816) (-1.3,-0.5624) (-1.25,-0.6484) (-1.2,-0.7369) (-1.15,-0.8243) (-1.1,-0.9068) (-1.05,-0.9805) (-1,-1.0420) (-0.95,-1.0880) (-0.9,-1.1165) (-0.85,-1.1262) (-0.8,-1.1169) (-0.75,-1.0892) (-0.7,-1.0447) (-0.65,-0.9858) (-0.6,-0.9151) (-0.55,-0.8356) (-0.5,-0.7501) (-0.45,-0.6614) (-0.4,-0.5721) (-0.35,-0.4843) (-0.3,-0.3995) (-0.25,-0.3191) (-0.2,-0.2439) (-0.15,-0.1743) (-0.1,-0.1106) (-0.05,-0.0526) (0,0.0000) (0.05,0.0476) (0.1,0.0907) (0.15,0.1300) (0.2,0.1661) (0.25,0.1997) (0.3,0.2314) (0.35,0.2620) (0.4,0.2921) (0.45,0.3225) (0.5,0.3540) (0.55,0.3874) (0.6,0.4237) (0.65,0.4641) (0.7,0.5098) (0.75,0.5627) (0.8,0.6247) (0.85,0.6987) (0.9,0.7883) (0.95,0.8983) (1,1.0355) (1.05,1.2090) (1.1,1.4320) (1.15,1.7232) (1.2,2.1099) (1.25,2.6329) (1.3,3.3540) (1.35,4.3686) (1.4,5.8275) (1.45,7.9736) (1.5,11.2080)};
\addplot[bbered, domain=-1.5:1.5] {x - x^2 + 2/3*x^3 + 2/3*x^4 - 16/15*x^5};
\legend{exacta, grado 5}
\end{axis}
\end{tikzpicture}
\caption{$y_2$ ($a_0=0,\ a_1=1$).}
\end{subfigure}
\caption{Ejemplo~\ref{ex:s2-tres}: solución exacta (calculada con la integral primera) y su polinomio de grado $5$.}
\label{fig:s2-ej3}
\end{figure}
 
El radio de convergencia no siempre es infinito aunque los coeficientes sean polinomios. El siguiente ejemplo muestra el fenómeno que ya vimos con $1/(1+x^2)$.
 
\begin{example}\label{ex:s2-raiz}
Resuelva en potencias de $x$ la ecuación $(1+x^2)y''+xy'-y=0$ e indique el radio de convergencia.
 
\noindent \textbf{Solución:} \quad En forma estándar, $P(x)=\dfrac{x}{1+x^2}$ y $Q(x)=-\dfrac{1}{1+x^2}$ son analíticas en $x=0$ y sus singularidades son $x=\pm i$, a distancia $1$ del origen: el Teorema~\ref{thm:s2-ordinario} garantiza $R\ge 1$. Sustituimos $y=\sum a_nx^n$ en la forma sin denominadores:
\[ \sum_{n=2}^{\infty} n(n-1)a_nx^{n-2}+\sum_{n=2}^{\infty} n(n-1)a_nx^{n}+\sum_{n=1}^{\infty} n a_nx^{n}-\sum_{n=0}^{\infty} a_nx^{n}=0. \]
Con $k$ como exponente común y agrupando $n(n-1)+n-1=n^2-1$,
\[ \sum_{k=0}^{\infty}\bigl[(k+2)(k+1)a_{k+2}+(k^2-1)a_k\bigr]x^k=0 \quad\Longrightarrow\quad a_{k+2}=-\frac{k-1}{k+2}\,a_k. \]
\begin{itemize}
\item \emph{Coeficientes impares:} $a_3=-\frac{0}{3}a_1=0$ y por lo tanto $a_5=a_7=\dots=0$. La solución impar es $y_2(x)=x$, ¡un polinomio! (Verificación: $0+x-x=0$.)
\item \emph{Coeficientes pares:} $a_2=\frac12a_0$, $a_4=-\frac14a_2=-\frac18a_0$, $a_6=-\frac36a_4=\frac1{16}a_0$, $a_8=-\frac58a_6=-\frac{5}{128}a_0$. Es la serie binomial de
\[ y_1(x)=\sqrt{1+x^2}=1+\tfrac12x^2-\tfrac18x^4+\tfrac1{16}x^6-\tfrac{5}{128}x^8+\cdots \]
\end{itemize}
La solución general es $y=a_0\sqrt{1+x^2}+a_1x$. El radio de $y_1$ es exactamente $R=1$ (razón: $\bigl|a_{k+2}/a_k\bigr|\to 1$ en $u=x^2$), aunque $y_1$ es suave en toda la recta real. \hfill $\square$
\end{example}
 
\begin{figure}[H]
\centering
\begin{tikzpicture}
\begin{axis}[bbeserie, height=6cm, xmin=-1.7, xmax=1.7, ymin=-0.7, ymax=3.2, restrict y to domain=-5:8,
  legend style={at={(0.5,-0.06)}, anchor=north}, samples=150, legend columns=3]
\addplot[black, very thick, domain=-1.7:1.7] {sqrt(1+x^2)};
\addplot[bbered, domain=-1.7:1.7] {1 + x^2/2};
\addplot[bbeamber, domain=-1.7:1.7] {1 + x^2/2 - x^4/8};
\addplot[bbegreen, domain=-1.7:1.7] {1 + x^2/2 - x^4/8 + x^6/16};
\addplot[bbeblue, domain=-1.7:1.7] {1 + x^2/2 - x^4/8 + x^6/16 - 5*x^8/128};
\legend{$\sqrt{1+x^2}$, hasta $x^2$, hasta $x^4$, hasta $x^6$, hasta $x^8$}
\end{axis}
\end{tikzpicture}
\caption{Ejemplo~\ref{ex:s2-raiz}: las sumas parciales convergen en $|x|<1$ y se alejan de la solución para $|x|>1$.}
\label{fig:s2-raiz}
\end{figure}
 
Un último ejemplo, muy importante en física, tiene una recurrencia que salta de \emph{tres} en tres.
 
\begin{example}\label{ex:s2-airy}
\textbf{(Ecuación de Airy.)} Determine dos soluciones en serie de $y''-xy=0$ centradas en $x=0$.
 
\noindent \textbf{Solución:} \quad $P=0$ y $Q=-x$ son polinomios: $R=\infty$. Con $y=\sum a_nx^n$,
\[ \sum_{k=0}^{\infty}(k+2)(k+1)a_{k+2}x^k-\sum_{k=1}^{\infty}a_{k-1}x^k=0 . \]
El término $k=0$ da $2a_2=0$, y para $k\ge 1$: $(k+2)(k+1)a_{k+2}=a_{k-1}$, es decir,
\[ a_2=0,\qquad a_{k+3}=\frac{a_k}{(k+3)(k+2)},\qquad k=0,1,2,\dots \]
La recurrencia conecta $a_0\to a_3\to a_6\to\cdots$ y $a_1\to a_4\to a_7\to\cdots$, mientras que $a_2=a_5=a_8=\cdots=0$:
\begin{center}
\small
\begin{tabular}{|c|c|c|c|c|}
\hline
\rowcolor{bbe!15}$m$ & $1$ & $2$ & $3$ & $4$ \\ \hline
$a_{3m}/a_0$ & $\frac{1}{2\cdot 3}=\frac1{6}$ & $\frac{1}{6\cdot 5\cdot 6}=\frac{1}{180}$ & $\frac{1}{180\cdot 8\cdot 9}=\frac{1}{12960}$ & $\frac{1}{12960\cdot 11\cdot 12}$ \\ \hline
$a_{3m+1}/a_1$ & $\frac{1}{3\cdot 4}=\frac1{12}$ & $\frac{1}{12\cdot 6\cdot 7}=\frac{1}{504}$ & $\frac{1}{504\cdot 9\cdot 10}=\frac{1}{45360}$ & $\frac{1}{45360\cdot 12\cdot 13}$ \\ \hline
\end{tabular}
\end{center}
Por lo tanto, dos soluciones linealmente independientes en $(-\infty,\infty)$ son
\[ y_1(x)=1+\frac{x^3}{6}+\frac{x^6}{180}+\frac{x^9}{12960}+\cdots,\qquad y_2(x)=x+\frac{x^4}{12}+\frac{x^7}{504}+\frac{x^{10}}{45360}+\cdots \tag*{$\square$} \]
\end{example}
 
\begin{figure}[H]
\centering
\begin{tikzpicture}
\begin{axis}[bbeosc, height=6.2cm, xmin=-8, xmax=2.5, ymin=-3.5, ymax=6, restrict y to domain=-8:12,
  legend pos=north west, samples=100]
\bbezero{-8}{2.5}
\addplot[bbeblue, very thick, smooth] coordinates {(-8,-0.3436) (-7.9,-0.2127) (-7.8,-0.0652) (-7.7,0.0874) (-7.6,0.2334) (-7.5,0.3617) (-7.4,0.4631) (-7.3,0.5305) (-7.2,0.5594) (-7.1,0.5482) (-7,0.4984) (-6.9,0.4138) (-6.8,0.3009) (-6.7,0.1676) (-6.6,0.0231) (-6.5,-0.1230) (-6.4,-0.2611) (-6.3,-0.3826) (-6.2,-0.4802) (-6.1,-0.5481) (-6,-0.5828) (-5.9,-0.5827) (-5.8,-0.5484) (-5.7,-0.4824) (-5.6,-0.3891) (-5.5,-0.2740) (-5.4,-0.1439) (-5.3,-0.0061) (-5.2,0.1321) (-5.1,0.2635) (-5,0.3815) (-4.9,0.4805) (-4.8,0.5561) (-4.7,0.6052) (-4.6,0.6259) (-4.5,0.6179) (-4.4,0.5822) (-4.3,0.5210) (-4.2,0.4374) (-4.1,0.3356) (-4,0.2200) (-3.9,0.0956) (-3.8,-0.0325) (-3.7,-0.1595) (-3.6,-0.2805) (-3.5,-0.3915) (-3.4,-0.4889) (-3.3,-0.5697) (-3.2,-0.6317) (-3.1,-0.6736) (-3,-0.6947) (-2.9,-0.6950) (-2.8,-0.6752) (-2.7,-0.6365) (-2.6,-0.5807) (-2.5,-0.5098) (-2.4,-0.4262) (-2.3,-0.3323) (-2.2,-0.2308) (-2.1,-0.1242) (-2,-0.0150) (-1.9,0.0946) (-1.8,0.2023) (-1.7,0.3064) (-1.6,0.4054) (-1.5,0.4979) (-1.4,0.5829) (-1.3,0.6598) (-1.2,0.7282) (-1.1,0.7878) (-1,0.8388) (-0.9,0.8814) (-0.8,0.9161) (-0.7,0.9435) (-0.6,0.9643) (-0.5,0.9793) (-0.4,0.9894) (-0.3,0.9955) (-0.2,0.9987) (-0.1,0.9998) (0,1.0000) (0.1,1.0002) (0.2,1.0013) (0.3,1.0045) (0.4,1.0107) (0.5,1.0209) (0.6,1.0363) (0.7,1.0578) (0.8,1.0868) (0.9,1.1245) (1,1.1723) (1.1,1.2319) (1.2,1.3050) (1.3,1.3938) (1.4,1.5008) (1.5,1.6288) (1.6,1.7813) (1.7,1.9624) (1.8,2.1770) (1.9,2.4308) (2,2.7309) (2.1,3.0858) (2.2,3.5056) (2.3,4.0028) (2.4,4.5923) (2.5,5.2924)};
\addplot[bbered, very thick, smooth] coordinates {(-8,-0.2676) (-7.9,-0.4530) (-7.8,-0.6027) (-7.7,-0.7058) (-7.6,-0.7550) (-7.5,-0.7471) (-7.4,-0.6835) (-7.3,-0.5696) (-7.2,-0.4144) (-7.1,-0.2295) (-7,-0.0284) (-6.9,0.1748) (-6.8,0.3660) (-6.7,0.5324) (-6.6,0.6635) (-6.5,0.7509) (-6.4,0.7899) (-6.3,0.7786) (-6.2,0.7184) (-6.1,0.6140) (-6,0.4722) (-5.9,0.3023) (-5.8,0.1145) (-5.7,-0.0798) (-5.6,-0.2697) (-5.5,-0.4446) (-5.4,-0.5952) (-5.3,-0.7138) (-5.2,-0.7947) (-5.1,-0.8345) (-5,-0.8319) (-4.9,-0.7880) (-4.8,-0.7055) (-4.7,-0.5893) (-4.6,-0.4455) (-4.5,-0.2812) (-4.4,-0.1043) (-4.3,0.0771) (-4.2,0.2554) (-4.1,0.4229) (-4,0.5732) (-3.9,0.7007) (-3.8,0.8010) (-3.7,0.8709) (-3.6,0.9087) (-3.5,0.9139) (-3.4,0.8872) (-3.3,0.8304) (-3.2,0.7463) (-3.1,0.6384) (-3,0.5106) (-2.9,0.3676) (-2.8,0.2140) (-2.7,0.0543) (-2.6,-0.1069) (-2.5,-0.2653) (-2.4,-0.4171) (-2.3,-0.5590) (-2.2,-0.6881) (-2.1,-0.8020) (-2,-0.8992) (-1.9,-0.9784) (-1.8,-1.0391) (-1.7,-1.0811) (-1.6,-1.1048) (-1.5,-1.1108) (-1.4,-1.1002) (-1.3,-1.0741) (-1.2,-1.0342) (-1.1,-0.9818) (-1,-0.9186) (-0.9,-0.8463) (-0.8,-0.7663) (-0.7,-0.6802) (-0.6,-0.5893) (-0.5,-0.4948) (-0.4,-0.3979) (-0.3,-0.2993) (-0.2,-0.1999) (-0.1,-0.1000) (0,0.0000) (0.1,0.1000) (0.2,0.2001) (0.3,0.3007) (0.4,0.4021) (0.5,0.5052) (0.6,0.6109) (0.7,0.7202) (0.8,0.8346) (0.9,0.9556) (1,1.0853) (1.1,1.2259) (1.2,1.3800) (1.3,1.5508) (1.4,1.7417) (1.5,1.9571) (1.6,2.2019) (1.7,2.4820) (1.8,2.8044) (1.9,3.1775) (2,3.6111) (2.1,4.1171) (2.2,4.7098) (2.3,5.4064) (2.4,6.2277) (2.5,7.1990)};
\legend{$y_1$, $y_2$}
\end{axis}
\end{tikzpicture}
\caption{Ejemplo~\ref{ex:s2-airy}: para $x<0$ las soluciones oscilan (con amplitud decreciente y frecuencia creciente); para $x>0$ crecen sin oscilar. La serie lo ``sabe'' todo: es válida en toda la recta.}
\label{fig:s2-airy}
\end{figure}
 
\begin{table}[H]
\centering
\small
\caption{Resumen de los ejemplos: la forma de la recurrencia sugiere cuánto se puede resolver.}
\label{tab:s2-resumen}
\begin{tabular}{|p{3.9cm}|p{4.2cm}|p{4.1cm}|c|}
\hline
\rowcolor{bbe!15}\textbf{Ecuación} & \textbf{Recurrencia} & \textbf{Solución cerrada} & $\mathbf{R}$ \\ \hline
$y''+y=0$ & $a_{k+2}=-\dfrac{a_k}{(k+2)(k+1)}$ & $\cos x,\ \sin x$ & $\infty$ \\ \hline
$y''-2xy'-4y=0$ & $a_{k+2}=\dfrac{2a_k}{k+1}$ & $y_2=xe^{x^2}$ & $\infty$ \\ \hline
$y''+(2-4x^2)y'-8xy=0$ & $a_{k+2}=\dfrac{4a_{k-1}-2a_{k+1}}{k+2}$ (3 términos) & vía $y'+(2-4x^2)y=C$ & $\infty$ \\ \hline
$(1+x^2)y''+xy'-y=0$ & $a_{k+2}=-\dfrac{k-1}{k+2}a_k$ & $x,\ \sqrt{1+x^2}$ & $1$ \\ \hline
$y''-xy=0$ & $a_{k+3}=\dfrac{a_k}{(k+3)(k+2)}$ (salto de 3) & funciones de Airy & $\infty$ \\ \hline
\end{tabular}
\end{table}
 
\begin{exercise}[Para llevar]
\begin{enumerate}
\item Resuelve $y''+xy'+y=0$ en potencias de $x$ e identifica una solución en términos de funciones elementales.
\item ¿Qué radio de convergencia garantiza el teorema para las soluciones de $(x^2-2x+10)y''+y'+y=0$ en potencias de $x$? ¿Y en potencias de $(x-1)$?
\end{enumerate}
\emph{Respuestas:} (1) $a_{k+2}=-\dfrac{a_k}{k+2}$, de donde $a_{2n}=\dfrac{(-1)^n}{2^n\,n!}a_0$ y $y_1=e^{-x^2/2}$; la otra solución es $y_2=\sum_{n\ge0}\dfrac{(-1)^n}{1\cdot3\cdots(2n+1)}x^{2n+1}$ (ambas con $R=\infty$). \; (2) Las raíces de $x^2-2x+10=(x-1)^2+9$ son $1\pm3i$: en potencias de $x$, $R\ge\sqrt{10}$; en potencias de $(x-1)$, $R\ge 3$.
\end{exercise}
